Every ten years, following the decennial census, the United States undergoes congressional redistricting. In most states, the legislature draws the lines. Elected officials design the districts they will run in. The result is a structural conflict of interest that has produced increasingly sophisticated partisan maps, from North Carolina's ``ugly gerrymander'' struck down in \textit{Rucho v. Common Cause}, 588 U.S.\ 684 (2019), to Maryland's ``Democratic gerrymander'' of the 6th Congressional District litigated in \textit{Benisek v. Lamone}.

Yet reform is rare. Between 2000 and 2024, forty-seven redistricting reform proposals were introduced across state legislatures. Thirty-nine failed in committee or on the floor. Of the eight that passed, five were substantially diluted before enactment. Why does an overwhelmingly popular reform---surveys consistently show 70--80\% of voters oppose partisan gerrymandering \citep{gallup2019}---fail so reliably through ordinary legislative channels?

The answer is straightforward but consequential: the legislators who must enact reform are the same officials who benefit from the status quo. This is not merely a matter of self-interest---though it is that---but a structural feature of any reform that requires the beneficiaries of a rule to vote for its abolition. Reform through the legislature asks incumbents to reduce the value of their most durable asset: the district lines that help ensure their re-election.

This paper develops a formal political economy model of redistricting reform. Section~2 presents the model, which treats favorable district lines as a capital asset yielding an incumbency rent each election cycle and shows why incumbents have strong incentives to block reform through ordinary channels. Section~3 identifies three bypass channels---direct democracy, judicial intervention, and federal legislation---that route around incumbent resistance. Section~4 tests the model's predictions against historical reform outcomes across 22 states. Section~5 discusses implications for the remainder of Track P.

The analysis has direct practical import. The \texttt{bisect} algorithmic redistricting platform described throughout this research program provides a technically sound and legally durable solution to partisan gerrymandering. But a technically sound solution does not automatically translate into adoption. Understanding why reform is hard, and which pathways have succeeded, is prerequisite to effective reform strategy. Papers P.1 through P.5 each develop one bypass channel in depth; this paper provides the analytical foundation for all five.
